\begin{proposition}[3.2.3]
\label{IV.2.3.2.3-fr}
Sean $X$ un preesquema localmente noetheriano, $U$ un abierto de $X$ y
$\mathcal F$ un $\mathcal O_X$-módulo coherente tal que
$U\cap\operatorname{Supp}(\mathcal F)$ sea denso en
$\operatorname{Supp}(\mathcal F)$. Las condiciones siguientes son
equivalentes:
\begin{enumerate}
\item[\emph{a)}] $\mathcal F$ es reducido.
\item[\emph{b)}] $\mathcal F|U$ es reducido y $\mathcal F$ carece de ciclo
primo asociado inmerso.
\item[\emph{c)}] Existen un subpreesquema cerrado reducido $X'$ de $X$ y un
$\mathcal O_{X'}$-módulo coherente $\mathcal F'$ sin torsión y de rango $1$
sobre toda componente irreducible de $X'$, tales que, si $j:X'\to X$ es la
inyección canónica, se tiene $j_*(\mathcal F')=\mathcal F$.
\end{enumerate}
Además, cuando esto ocurre, el subpreesquema $X'$ está definido por el ideal
$\mathcal J$ de $\mathcal O_X$ anulador de $\mathcal F$.
\end{proposition}

Para ver que \emph{c)} implica \emph{a)}, es posible, en virtud de (3.1.3),
limitarse al caso en que $X'=X$ es íntegro, de punto genérico $x$, y se puede
suponer además $X=\operatorname{Spec}(A)$ afín y $\mathcal F=\widetilde M$,
donde $A$ es, por tanto, íntegro y $M$ es un $A$-módulo sin torsión de rango
$1$; como el anulador de todo elemento de $M$ se reduce entonces a $0$, se
tiene $\operatorname{Ass}(\mathcal F)=\{x\}$ y $\mathcal F_x$ es isomorfo a
$\mathcal O_x$, cuerpo de fracciones de $A$, por lo que se satisfacen las
condiciones de la definición (3.2.2). Como la existencia de ciclos primos
asociados inmersos es local, es claro que, si $\mathcal F$ no tiene tales
ciclos y $U\cap\operatorname{Supp}(\mathcal F)$ es denso en
$\operatorname{Supp}(\mathcal F)$, se tiene
$\operatorname{Ass}(\mathcal F)=\operatorname{Ass}(\mathcal F|U)$, y por
tanto \emph{a)} y \emph{b)} son equivalentes. Si se satisface \emph{a)},
tomemos como $X'$ el subpreesquema cerrado reducido de $X$ cuyo espacio
subyacente es $\operatorname{Supp}(\mathcal F)$ (\textbf I, 5.2.1), y sea
$\mathcal F'=j^*(\mathcal F)$; un punto $x$ de
$\operatorname{Ass}(\mathcal F)$ es necesariamente un punto maximal de $X'$ y,
como $\operatorname{long}(\mathcal F_x)=1$, $\mathcal F'_x$ es isomorfo a
$\mathcal F_x$, y por tanto al cuerpo $\boldsymbol{k}(x)=\mathcal O_{X',x}$, lo que prueba
que $\mathcal F'$ es sin torsión y de rango $1$ (\textbf I, 7.4.6 y 7.4.2). Por
último, la última afirmación es trivial, puesto que para todo $y\in X'$ el
anulador del $\mathcal O_{X',y}$-módulo $\mathcal F'_y$ es nulo.

